\documentclass[a4paper]{article}
% Español (México) con XeLaTeX: fontspec para fuentes Unicode, polyglossia
% para la división silábica y los nombres del documento en la variante mexicana.
\usepackage{fontspec}
\usepackage{polyglossia}
\setdefaultlanguage[variant=mexican]{spanish}
% Los nombres chinos van como «pinyin (original)»: xeCJK da a los caracteres
% Han su propia fuente; sin ella XeLaTeX los omite con un aviso.
% FandolSong no cubre algunos caracteres de nombres (U+73FA); WenQuanYi Zen Hei sí.
\usepackage[AutoFallBack=true]{xeCJK}
\setCJKmainfont{FandolSong-Regular.otf}
\setCJKfallbackfamilyfont{\CJKrmdefault}{WenQuanYi Zen Hei}
% polyglossia no traduce los nombres de listings.
\AtBeginDocument{\providecommand{\lstlistingname}{}\renewcommand{\lstlistingname}{Listado}%
  \providecommand{\lstlistlistingname}{}\renewcommand{\lstlistlistingname}{Índice de listados}}
\usepackage{amsmath, amssymb}
\usepackage{graphicx}
\usepackage[margin=2.35cm]{geometry}
\usepackage[most]{tcolorbox}
\usepackage{booktabs}
\usepackage{tabularx}
\usepackage{longtable}
\usepackage{array}
\usepackage{float}
\usepackage{xcolor}
\usepackage{hyperref}
\usepackage{fancyhdr}
\hypersetup{colorlinks=true, linkcolor=blue!55!black, urlcolor=blue!60!black}
\graphicspath{{./}{figures/}}
\setlength{\parskip}{0.55em}
\setlength{\parindent}{2em}
\setlength{\emergencystretch}{3em}
\renewcommand{\arraystretch}{1.18}
\newtcolorbox{knowledgebox}[1]{enhanced, colback=blue!5!white, colframe=blue!70!black, colbacktitle=blue!70!black, coltitle=white, fonttitle=\bfseries, title={#1}, sharp corners, breakable}
\newtcolorbox{importantbox}[1]{enhanced, colback=yellow!10!white, colframe=yellow!75!black, colbacktitle=yellow!75!black, coltitle=black, fonttitle=\bfseries, title={#1}, sharp corners, breakable}
\newtcolorbox{warningbox}[1]{enhanced, colback=red!5!white, colframe=red!70!black, colbacktitle=red!70!black, coltitle=white, fonttitle=\bfseries, title={#1}, sharp corners, breakable}
\pagestyle{fancy}
\fancyhf{}
\fancyfoot[C]{\thepage}
\fancyfoot[R]{\footnotesize\color{gray}\href{https://github.com/hqhq1025/ai-course-notes}{AI Course Notes}}
\renewcommand{\headrulewidth}{0pt}
\renewcommand{\footrulewidth}{0pt}
\newcommand{\videourl}{https://www.youtube.com/watch?v=9zSMTUUEfmU}
\newcommand{\repourl}{https://github.com/hqhq1025/ai-course-notes}

\begin{document}
\begin{titlepage}
\centering
{\Large Notas didácticas de una entrevista a fondo\par}
\vspace{1.0cm}
{\huge\bfseries Zhipu (智谱): de la inteligencia cognitiva a la primera empresa de grandes modelos en cotizar en bolsa en el mundo\par}
\vspace{0.35cm}
{\Large Zhang Xiaojun conversa con Zhang Peng (张鹏)\par}
\vspace{0.35cm}
{\large Elaborado a partir del video público de Zhang Xiaojun Podcast, la transcripción hecha en GPU, la estructura de capítulos y diagramas conceptuales propios\par}
\vspace{0.3cm}
{\large 2026-01-08\par}
\vspace{0.85cm}
\includegraphics[width=0.88\textwidth,height=0.42\textheight,keepaspectratio]{cover.jpg}\par
\vfill
\begin{tcolorbox}[width=0.92\textwidth, colback=black!2!white, colframe=black!55, sharp corners]
\textbf{Título del video}: 129. Entrevista sobre la salida a bolsa de la primera empresa de grandes modelos en cotizar en el mundo; conversación con Zhang Peng, CEO de Zhipu: ¿hacia dónde va el camino?\par
\textbf{Canal del video}: Zhang Xiaojun Podcast\par
\textbf{Duración del video}: 02:26:40\par
\textbf{Enlace del video}: \href{\videourl}{\nolinkurl{\videourl}}\par
\textbf{Nota sobre los materiales}: El video de YouTube no tiene subtítulos disponibles; el texto principal se elaboró a partir de la transcripción con faster-whisper large-v3 en CUDA, los capítulos del video, la portada y figuras didácticas propias. Las tomas fijas de la entrevista no se reutilizan en el texto principal.\par
\textbf{Página del proyecto}: \href{\repourl}{\nolinkurl{\repourl}}\par
\end{tcolorbox}
\end{titlepage}

\tableofcontents
\newpage
\section{Introducción: salir a bolsa no es el punto final, sino una revisión del camino recorrido}

Esta entrevista ocurrió en un momento particular: Zhipu (智谱) confirmó que el 8 de enero de 2026 empezaría a cotizar en la Bolsa de Hong Kong, con lo que se convirtió en la primera empresa china de modelos grandes en salir a bolsa; también se le ha llamado «la primera acción de modelos grandes del mundo». Zhang Peng (张鹏) dio la entrevista con una pierna lesionada por una caída y apoyado en muletas, y el programa usa «Break a leg» como metáfora de apertura. Pero lo verdaderamente importante de este episodio no es el dramatismo del toque de campana en la bolsa, sino cómo Zhipu, a lo largo de más de seis años, pasó de la exploración de la inteligencia cognitiva en un laboratorio de Tsinghua a la industria de los modelos grandes y al gobierno corporativo de una empresa que cotiza en bolsa.

Estas notas lo organizan en tres líneas. La primera es la organizacional: cómo la tradición P2P del Laboratorio de Ingeniería del Conocimiento de Tsinghua se convirtió en la vía de transferencia de resultados de investigación de Zhipu. La segunda es la tecnológica: de la inteligencia perceptiva a la inteligencia cognitiva, sacudida una y otra vez por GPT-3, ChatGPT, DeepSeek y la Scaling Law. La tercera es la empresarial: cómo una empresa de IA de origen académico maneja el código abierto frente al código cerrado, el idealismo frente a la comercialización y el objetivo de largo plazo de la AGI frente a la disciplina de corto plazo de una empresa que cotiza en bolsa.

\begin{importantbox}{Tesis central de este episodio}
La historia de Zhipu no es la de «alcanzar de repente la ola de los modelos grandes», sino la de haber mantenido durante mucho tiempo la investigación, la ingeniería, la industria y la comercialización en una misma cadena. Salir a bolsa es solo un resultado de etapa; Zhang Peng preferiría que la historia registrara a Zhipu como «precursora de la AGI» y como «la que abrió el camino».
\end{importantbox}

\begin{knowledgebox}{Nota sobre la estrategia visual}
Este video es una toma fija de entrevista, sin slides, pizarrón ni demostraciones de producto. El cuerpo del texto no repite fotogramas de las personas; la portada sirve para identificar la fuente, y en el cuerpo el contenido didáctico se apoya en diagramas conceptuales sobre la transferencia de resultados, los sacudimientos causados por los modelos, la Scaling Law, el código abierto frente al cerrado y las etapas de la empresa.
\end{knowledgebox}

\subsection{Resumen de la sección}

Conviene leer EP129 junto con el informe trimestral sobre modelos grandes de EP136 y la historia de la tecnología de Agent de EP139. Desde la perspectiva interna de una empresa de modelos, explica cómo la industria de los modelos grandes toma forma por la acción conjunta de los paradigmas tecnológicos, los mercados de capitales, el ecosistema de código abierto y el gobierno de las organizaciones.
\section{Los pioneros: del laboratorio de Tsinghua a la transferencia de resultados científicos y tecnológicos}

Al inicio de la entrevista, Zhang Peng (张鹏) regresa a los orígenes de Zhipu (智谱). Zhipu no surgió de la nada como emprendimiento: creció a partir de la investigación y de la tradición de transferencia a la ingeniería del Laboratorio de Ingeniería del Conocimiento del Departamento de Ciencias de la Computación de la Universidad Tsinghua (清华大学). La generación anterior de los «cuatro pequeños dragones» de la inteligencia artificial le mostró al laboratorio que la industrialización de la IA era posible, pero también le hizo ver que, en la etapa de la inteligencia perceptiva, el techo de capacidades era evidente y que la siguiente generación de inteligencia artificial necesitaba una nueva dirección.

\begin{figure}[H]
\centering
\includegraphics[width=0.96\textwidth]{zhipu-ai-roadmap.png}
\caption{La trayectoria de seis años de Zhipu: de la inteligencia cognitiva y la transferencia P2P hasta la salida a bolsa, y de regreso a la meta de la AGI.}
\end{figure}

\begin{knowledgebox}{Lectura de la figura: esta no es una trayectoria puramente técnica}
Cada nodo de la figura tiene a la vez un significado técnico y uno organizacional. La inteligencia cognitiva es una dirección técnica; P2P es un mecanismo de transferencia de resultados; GPT-3 y ChatGPT son choques de paradigma externos; DeepSeek cambia el discurso sobre el código abierto y la eficiencia; la IPO, por su parte, lleva a la empresa a una nueva etapa de gobierno corporativo.
\end{knowledgebox}

\subsection{P2P: Paper to Product}

Lo anterior describe la trayectoria general de Zhipu; aquí se entra en su método organizacional más básico. P2P explica por qué Zhipu coloca de forma natural los artículos, los sistemas, los productos y los clientes en una misma cadena.

Zhang Peng menciona una y otra vez la tradición P2P de Tsinghua: paper to project / paper to product. Los resultados de investigación no pueden quedarse en artículos y prototype, sino que deben convertirse en system y product utilizables por clientes empresariales. Lo que él tenía a su cargo en el laboratorio era precisamente la transferencia a la ingeniería: los profesores y los estudiantes investigaban, el equipo de ingeniería convertía los resultados de investigación en sistemas y después atendía las necesidades de los clientes.

\begin{figure}[H]
\centering
\includegraphics[width=0.94\textwidth]{paper-to-product-loop.png}
\caption{P2P: del artículo al producto, y de vuelta a la investigación mediante la retroalimentación de los clientes.}
\end{figure}

\begin{knowledgebox}{Lectura de la figura: P2P es el material genético organizacional de Zhipu}
La figura va de Paper a Prototype, luego a System, Product y Feedback. Muestra que Zhipu no es un instituto de investigación puramente académico ni una empresa puramente comercial; sus capacidades iniciales provienen de llevar la investigación de frontera a la ingeniería y de dejar que las necesidades de clientes reales corrijan, a su vez, la investigación.
\end{knowledgebox}

\subsection{Los primeros en transferir resultados científicos y tecnológicos}

Si la sección anterior presentó P2P como una capacidad, esta examina la vía institucional. Si la transferencia de resultados científicos y tecnológicos no se concreta, aunque la investigación de frontera llegue a convertirse en un sistema, difícilmente se transformará en un activo empresarial capaz de financiarse, operar a largo plazo y cotizar en bolsa.

Hacia 2018, la política nacional abrió una ventana para que el personal de las instituciones de investigación transfiriera resultados científicos y tecnológicos, pero los detalles no estaban claros: cómo valuar los resultados, cómo repartir entre la universidad y el equipo, cómo cumplir con la normativa en el proceso y quién asumiría las responsabilidades posteriores. El equipo de Zhipu se convirtió en «los primeros en comer cangrejo»: negociaba con la universidad mientras exploraba la vía institucional, y solo después de más de un año logró concretar el registro de la empresa y la salida del equipo.

\begin{warningbox}{La transferencia de resultados no es simplemente «un profesor que abre una empresa»}
La transferencia de resultados universitarios involucra propiedad intelectual, activos estatales, los derechos de la universidad, los incentivos del equipo fundador y el gobierno corporativo. Si la vía institucional no es clara, el financiamiento posterior, la salida a bolsa y el cumplimiento normativo quedarán expuestos a riesgos. Que Zhipu eligiera la vía formal fue una condición importante para la permanencia de la empresa a largo plazo.
\end{warningbox}

\subsection{Por qué la vía institucional afecta a una empresa tecnológica}

Las empresas tecnológicas suelen concentrar su atención en los modelos, los artículos, el financiamiento y los productos, pero la historia temprana de Zhipu nos recuerda que la vía institucional también es una capacidad. Si en una empresa surgida de una universidad o de una institución de investigación no está clara la titularidad de la propiedad intelectual, la distribución de las acciones ni los derechos de la universidad, a corto plazo quizá no se noten problemas, pero a largo plazo, al llegar el financiamiento, las auditorías, la salida a bolsa y la cooperación internacional, eso se convierte en un riesgo estructural.

Zhipu dedicó más de un año a este asunto; parece lento, pero absorbió por adelantado los riesgos futuros. Esta decisión es coherente con el posterior «carácter de cemento» de la empresa: no busca la mayor notoriedad en el menor tiempo, sino asentar bien los cimientos. Para una empresa de IA, este discreto trabajo institucional llama menos la atención que el lanzamiento de un modelo, pero determina si la empresa podrá llegar lejos.

\begin{knowledgebox}{Las cuatro cuestiones de la transferencia de resultados científicos y tecnológicos}
\begin{tabularx}{\textwidth}{p{0.22\textwidth}p{0.34\textwidth}X}
\toprule
Cuestión & Qué debe resolver & Qué pasa si no se resuelve \\
\midrule
Titularidad de los resultados & A quién pertenecen los artículos, el código, las patentes y los sistemas & Surgen disputas de propiedad al financiarse y al salir a bolsa. \\
Valuación & Cuánto valor comercial representan los resultados de investigación & Es difícil justificar la distribución de acciones y el cumplimiento en materia de activos estatales. \\
Incentivos del equipo & Cómo se reparten los derechos entre el equipo fundador y la universidad & Falta motivación entre los primeros colaboradores o surgen disputas posteriores. \\
Vía de cumplimiento & Si se cumplen los requisitos de la universidad y de la política pública & Cuanto más crece la empresa, más difícil es subsanar los problemas heredados. \\
\bottomrule
\end{tabularx}
\end{knowledgebox}

\subsection{Resumen de la sección}

Los inicios de Zhipu no fueron un «emprendimiento de modelos», sino la ingeniería aplicada a la investigación y la transferencia de resultados científicos y tecnológicos. El material genético P2P explica por qué desde el principio dio importancia al ciclo cerrado entre la investigación, los sistemas, los clientes y la comercialización.
\section{Inteligencia cognitiva: el siguiente paso después de la inteligencia perceptiva}

Entre 2015 y 2016, la generación anterior de empresas de IA se centraba sobre todo en llevar a la práctica la CV, el NLP temprano y el aprendizaje profundo. Zhang Peng (张鹏) cuenta que en ese momento se dieron cuenta de que la inteligencia perceptiva tenía un techo: una vez que el reconocimiento facial superó al ser humano, seguir mejorándolo tenía un sentido limitado, y su efecto en la industria también se acercaba a su límite. Por eso propusieron la «inteligencia cognitiva» como el siguiente escalón hacia la inteligencia artificial general.

\begin{figure}[H]
\centering
\includegraphics[width=0.94\textwidth]{cognitive-intelligence-ladder.png}
\caption{De la inteligencia perceptiva a la inteligencia cognitiva, y de ahí a los grandes modelos, los Agent y la AGI.}
\end{figure}

\begin{knowledgebox}{Lectura de la figura: la inteligencia cognitiva es una definición por etapas}
En la figura, la inteligencia cognitiva no es la AGI en sí, sino el siguiente paso después de la inteligencia perceptiva. Busca llevar la IA de las tareas de reconocimiento, clasificación y percepción a la comprensión, el razonamiento, el conocimiento, el lenguaje y las tareas complejas. Más tarde, los grandes modelos se convirtieron en la forma tecnológica principal de esa dirección.
\end{knowledgebox}

\subsection{Por qué no definir directamente la AGI}

Zhang Peng explica que en ese momento no se empeñaron en definir la AGI, porque estaba demasiado lejos y no podía definirse con claridad; en cambio, sí podía intentarse definir «cuál es el siguiente paso». Es un enfoque prudente: no sustituir un problema de la etapa actual por un concepto último. La inteligencia cognitiva le dio a Zhipu (智谱) una dirección lo bastante lejana, pero que aún podía avanzarse mediante la ingeniería.

\begin{importantbox}{El valor de los objetivos por etapas}
Una visión de largo plazo necesita definiciones por etapas. La AGI es un objetivo lejano; la inteligencia cognitiva es un objetivo de etapa que puede discutirse, que permite organizar recursos y que admite una conversión a ingeniería.
\end{importantbox}

\subsection{Resumen de la sección}

La inteligencia cognitiva fue el punto de anclaje de la dirección de Zhipu en sus inicios. Le permitió a la empresa no quedarse en la inteligencia perceptiva ni proclamar la AGI sin sustento, sino buscar una ruta hacia la siguiente generación de IA a partir de la comprensión, el razonamiento y la ingeniería del conocimiento.
\section{GPT-3 y ChatGPT: cómo los impactos externos reescribieron el criterio de la empresa}

La inteligencia cognitiva marcó la dirección interna, pero la ola de los grandes modelos reescribió una y otra vez esa ruta desde fuera. GPT-3 y ChatGPT fueron dos impactos de distinto nivel: uno cambió la imaginación técnica; el otro, a los usuarios y al mercado.

Que la entrevista trate a GPT-3 y a ChatGPT en capítulos separados es importante. GPT-3 representa la entrada del paradigma de escala en el horizonte; ChatGPT representa el momento en que la percepción de los usuarios y el ritmo de la industria se encendieron de golpe. Cuando Zhang Peng (张鹏) recuerda estos momentos, lo hace con un tono «a la vez ansioso y emocionado»: la ansiedad venía de que la ola externa era demasiado grande; la emoción, de que la dirección para la que se habían preparado durante tanto tiempo por fin quedaba validada.

\begin{figure}[H]
\centering
\includegraphics[width=0.96\textwidth]{model-shockwaves.png}
\caption{Tres impactos de modelos: GPT-3, ChatGPT, DeepSeek y los cambios de paradigma posteriores.}
\end{figure}

\begin{knowledgebox}{Lectura de la figura: cada impacto cambia un nivel}
GPT-3 cambió la imaginación técnica: la escala podría producir capacidades generales. ChatGPT cambió la percepción de los usuarios: cualquier persona supo lo que un gran modelo puede hacer. DeepSeek cambió el discurso sobre eficiencia y código abierto: la capacidad de un modelo no depende solo de invertir recursos masivos; también puede impulsarse con ingeniería y con un ecosistema abierto.
\end{knowledgebox}

\subsection{Por qué «la oportunidad es para quien está preparado»}

Después de los impactos externos, esta sección responde a un malentendido frecuente: si aprovechar una oportunidad depende solo de la suerte. La respuesta de Zhang Peng se acerca más a la acumulación cotidiana que a una predicción prodigiosa.

Zhang Peng cita a su mentor: una oportunidad es como un tablón en el mar; cuando llega flotando, todavía tienes que dar unas cuantas brazadas para alcanzarlo. La ola de los grandes modelos es muy difícil de predecir con precisión, pero sí es posible prepararse a largo plazo: mantener la dirección correcta y acumular sin pausa capacidades de ingeniería, de investigación y de comercialización, para tener con qué aprovechar la oportunidad cuando llegue.

\begin{importantbox}{Prepararse no es predecir}
Predecir es saber qué ocurrirá mañana; prepararse es seguir acumulando capacidades en la dirección correcta aunque no se sepa qué ocurrirá mañana. La ola de los grandes modelos se parece más a lo segundo.
\end{importantbox}

\subsection{Resumen de la sección}

La ruta de modelos de Zhipu (智谱) no consistió en seguir sin más a GPT-3 o a ChatGPT, sino en que su dirección temprana de inteligencia cognitiva fue corregida repetidamente por los paradigmas técnicos externos. Cuanto mayor es el impacto externo, más se pone a prueba si la organización está preparada.
\section{Evolución del paradigma de la Scaling Law}

La sección anterior trató el impacto de GPT y ChatGPT; esta entra en un paradigma técnico más fundamental. Al principio, la Scaling Law parecía un problema de escala en un solo eje; con el tiempo se convirtió en un problema de crecimiento de capacidades en varios ejes.

En la entrevista, «la evolución del paradigma de la Scaling Law» es uno de los ejes técnicos centrales. Zhang Peng (张鹏) señaló que el futuro no consiste solo en que siga mejorando la capacidad de los modelos fundacionales, sino también en la fusión de múltiples modelos y múltiples datos, los agentes y nuevas formas de cómputo; en particular, la línea de RL seguirá produciendo métodos nuevos. Es decir, la Scaling Law ya no equivale únicamente a la fórmula de preentrenamiento de «más parámetros, más datos, más cómputo».

\begin{figure}[H]
\centering
\includegraphics[width=0.96\textwidth]{scaling-law-evolution.png}
\caption{Evolución del paradigma de la Scaling Law: escala, multimodalidad, posentrenamiento, RL y Agent.}
\end{figure}

\begin{knowledgebox}{Lectura de la figura: la Scaling Law pasa de un eje a varios ejes}
El relato inicial de los grandes modelos ponía el acento en los parámetros, los datos y el cómputo; después se sumaron los datos multimodales, el posentrenamiento, la retroalimentación de preferencias y de tareas, el RL y los entornos de Agent. El crecimiento de capacidades no proviene solo de «ser más grande», sino también de «aprovechar mejor la retroalimentación».
\end{knowledgebox}

\subsection{Por qué el Agent se convirtió en la vía de aplicación}

Zhang Peng considera que los agentes son importantes porque resuelven el problema de cómo llevar un modelo hasta su aplicación real. El modelo es la base de capacidades, pero lo que el cliente quiere son resultados de negocio. El Agent conecta el modelo con herramientas, procesos, entornos y necesidades de los usuarios; es la capa intermedia para la comercialización y para liberar las capacidades.

\begin{knowledgebox}{Términos clave: Scaling Law y conceptos relacionados}
\begin{tabularx}{\textwidth}{p{0.23\textwidth}p{0.32\textwidth}X}
\toprule
Término & Explicación en una frase & Papel en este episodio \\
\midrule
Scaling Law & Regularidad con la que cambia la capacidad de un modelo al escalar parámetros, datos y cómputo & Relato de fondo del paradigma de los grandes modelos. \\
Post-training & Optimización continua del modelo después del preentrenamiento mediante preferencias, tareas y retroalimentación & Hace que el modelo sea más útil y más controlable. \\
RL & Reinforcement Learning, aprendizaje por refuerzo & Zhang Peng considera que traerá nuevas formas de cómputo. \\
Agent & Agente que invoca herramientas y ejecuta tareas & Vía de aplicación que conecta el modelo con usos reales. \\
Fusión de múltiples modelos & Varios modelos y fuentes de datos que en conjunto forman un sistema de capacidades & Una de las direcciones futuras de los modelos base. \\
\bottomrule
\end{tabularx}
\end{knowledgebox}

\subsection{Resumen de la sección}

La evolución de la Scaling Law muestra que la competencia entre grandes modelos ya no se reduce a «quién es más grande». La retroalimentación, el posentrenamiento, el Agent, el RL y los datos multimodales se están convirtiendo en nuevos ejes de crecimiento de capacidades.
\section{La diferenciación entre las empresas de modelos: la apuesta particular de Zhipu (智谱)}

En la entrevista, el conductor plantea una pregunta crucial: si OpenAI se parece cada vez más a una empresa de aplicaciones, Anthropic es fuerte en 2B y en coding, y Google, Meta, DeepSeek, Kimi, MiniMax, Qwen y otros ya ocupan cada uno su propia posición, ¿en qué se distingue Zhipu? La respuesta de Zhang Peng (张鹏) no es un producto puntual, sino tres predicciones y una apuesta de largo plazo.

Menciona tres predicciones que hizo a principios de 2025: primero, la capacidad de los modelos base seguirá mejorando, en especial la de los modelos base híbridos que integran múltiples modelos y diversos tipos de datos; segundo, los agentes inteligentes se convertirán en una dirección importante; tercero, la internacionalización ganará importancia. En retrospectiva, la evolución de la industria confirmó las tres. La palabra que define la apuesta de largo plazo es AGI; en el corto plazo, se desglosa en Agent y en nuevas formas de cómputo, sobre todo RL.

\subsection{Tres tipos de diferenciación entre empresas}

\begin{longtable}{p{0.23\textwidth}p{0.32\textwidth}p{0.36\textwidth}}
\toprule
Dirección de diferenciación & Ventaja típica & Pregunta que Zhipu debe responder \\
\midrule
Empresa de aplicaciones & Mejor comprensión del usuario, de la distribución y del ciclo completo de producto & Cómo logra Zhipu que el modelo llegue a aplicaciones reales, en lugar de ofrecer solo una API. \\
Empresa corporativa/2B & Mejor comprensión del cliente, del cumplimiento normativo, del servicio y de la entrega & Cómo convierte Zhipu su capacidad de investigación en ingresos sostenibles. \\
Plataforma de código abierto & Mayor difusión del ecosistema y más retroalimentación de los desarrolladores & Cómo transforma Zhipu su influencia en el código abierto en rendimientos compuestos, comerciales y técnicos. \\
Empresa de modelos base & Mayor capacidad del modelo y un sistema de entrenamiento más sólido & Cómo sostiene Zhipu la inversión y construye capacidades diferenciadas. \\
\bottomrule
\end{longtable}

\begin{importantbox}{La apuesta de Zhipu no es un solo punto}
Zhipu debe responder a la vez por la capacidad del modelo, la implementación de Agent, la internacionalización, el ecosistema de código abierto, la comercialización y el gobierno corporativo como empresa cotizada. No se parece a una empresa de aplicaciones de una sola línea ni a una comunidad de código abierto pura; se parece más a una empresa de modelos que mantiene el equilibrio en varias dimensiones.
\end{importantbox}

\subsection{Por qué Agent es el punto de apoyo en el corto plazo}

Zhang Peng afirma que Agent resuelve el problema de la ruta de implementación entre el modelo y la aplicación real. Esta afirmación es central. El modelo base es muy capaz por sí mismo, pero lo que el cliente paga es el resultado del trabajo; Agent conecta el modelo, las herramientas, los flujos de trabajo y los escenarios de negocio, de modo que las organizaciones puedan absorber la capacidad del modelo.

Desde esta perspectiva, EP129 y EP130 se complementan. Zhang Yueguang (张月光), desde la óptica del gerente de producto, sostiene que un Agent necesita ocupar un lugar en la mente del usuario y contar con retroalimentación de ciclo corto y alta frecuencia; Zhang Peng, desde la óptica de la empresa de modelos, sostiene que Agent es la capa intermedia por la que la capacidad del modelo entra en las aplicaciones. Ambos están en realidad en los dos lados del mismo problema: la empresa de modelos debe llevar la capacidad hasta el producto, y la empresa de productos debe organizar esa capacidad en un punto de entrada que los usuarios efectivamente utilicen.

\begin{knowledgebox}{Conexión con EP130}
EP130 explica cómo Dokie entra al Agent de creación y expresión a partir de PPT; EP129 explica cómo Zhipu concibe Agent como la ruta de implementación del modelo. El primero se ocupa del lugar en la mente del usuario; el segundo, de liberar la capacidad del modelo. Solo en conjunto forman el problema completo de la industrialización de Agent.
\end{knowledgebox}

\subsection{Resumen de la sección}

Las empresas de modelos ya comenzaron a diferenciarse. La posición particular de Zhipu consiste en que posee un ADN de investigación académica y, al mismo tiempo, debe resolver bien la comercialización, Agent, el código abierto y la internacionalización bajo el gobierno corporativo de una empresa cotizada.
\section{DeepSeek, código abierto y código cerrado}

La sección anterior sobre la Scaling Law trató el crecimiento de las capacidades; esta sección pasa a cómo la industria percibe esas capacidades y cómo se difunden. DeepSeek llevó el código abierto, la eficiencia y los resultados reales al centro del discurso de la industria.

En 2025, el título de esta sección fue «Aprender de DeepSeek». Zhang Peng (张鹏) mencionó que, en cierta medida, DeepSeek también le enseñó algo a la industria: en lo técnico no siempre hacen falta grandes eventos de lanzamiento, porque al final lo que cuenta es el resultado en aplicaciones reales. DeepSeek cambió a la vez el discurso sobre el código abierto y el discurso sobre la eficiencia, y llevó a las empresas de modelos a replantear la relación entre apertura, ecosistema y control comercial.

\begin{figure}[H]
\centering
\includegraphics[width=0.94\textwidth]{open-vs-closed.png}
\caption{Código abierto frente a código cerrado: el equilibrio entre la influencia en el ecosistema y el control comercial.}
\end{figure}

\begin{knowledgebox}{Lectura de la figura: el código abierto y el código cerrado no son una dicotomía moral}
El código abierto permite ampliar el ecosistema, recibir retroalimentación de los desarrolladores y construir una reputación técnica; el código cerrado permite conservar el control comercial, atender a clientes empresariales y manejar los límites de seguridad y cumplimiento normativo. Por lo general, las empresas de modelos necesitan combinar ambos enfoques en distintas capas de producto y distintos segmentos de clientes.
\end{knowledgebox}

\subsection{La retroalimentación global que trae el código abierto}

Zhang Peng mencionó que, después de que Zhipu (智谱) liberó sus modelos como código abierto, la evaluación en el extranjero fue muy positiva, e incluso hubo otros productos o empresas que usaron, destilaron o recortaron directamente los modelos. Esto muestra que el código abierto tiene influencia, pero también que da lugar a un ecosistema complejo en el que los modelos se reutilizan, se modifican y se vuelven a empaquetar.

\begin{warningbox}{El código abierto no solo trae elogios}
El código abierto aporta reputación entre los desarrolladores y difusión en el ecosistema, pero también trae problemas como la atribución comercial difusa, el reempaquetado de los modelos por terceros y la dificultad para capturar ingresos. La estrategia de código abierto debe diseñarse junto con la estrategia de comercialización.
\end{warningbox}

\subsection{Resumen de la sección}

Después de DeepSeek, la competencia entre empresas de modelos da más peso a la eficiencia, a los resultados reales y al ecosistema de código abierto. El código abierto y el código cerrado no son opciones excluyentes: la cuestión es cómo cada empresa de modelos reparte la influencia en el ecosistema, el valor comercial y los límites de seguridad.
\section{De la presentación de lanzamiento al desempeño real: cambios en la visibilidad de las empresas de modelos}

En la segunda mitad, Zhang Peng (张鹏) menciona un cambio muy sutil. Antes, el lanzamiento de un modelo solía requerir una presentación solemne, cientos de asistentes y una fuerte campaña de difusión. En cambio, tras la liberación en código abierto de GLM-4.5, Zhipu (智谱) no hizo un lanzamiento especialmente grande y, aun así, desarrolladores y productos del extranjero lo valoraron muy bien. Incluso hubo productos como Windsurf que sustituyeron su modelo o integraron este, y otros fabricantes que lo destilaron o recortaron. Este cambio indica que la visibilidad técnica de las empresas de modelos grandes está pasando de la «narrativa de la presentación de lanzamiento» al «desempeño en el uso real».

\subsection{Por qué disminuye el rendimiento marginal de las presentaciones de lanzamiento}

Antes se dijo que la visibilidad de un modelo se desplaza hacia su desempeño real; esta sección explica por qué. Cuando se lanzan demasiados modelos, la narrativa deja de ser escasa; lo escaso pasa a ser el uso real y la eficiencia en costos.

Después de varias rondas de lanzamientos en la industria de los modelos grandes, usuarios y desarrolladores ya se cansaron de «otro lanzamiento de modelo más». La publicidad puede atraer atención a corto plazo, pero que un modelo llegue a integrarse en el trabajo de los desarrolladores, que los clientes paguen por él y que la comunidad lo discuta por iniciativa propia depende cada vez más de su capacidad real y de su eficiencia en costos. La influencia de DeepSeek también radica en que llevó al primer plano de la industria el camino de «hablar poco, liberar mucho en código abierto y dejar que los resultados hablen».

\begin{knowledgebox}{Tres fuentes de visibilidad de un modelo}
\begin{tabularx}{\textwidth}{p{0.22\textwidth}p{0.34\textwidth}X}
\toprule
Fuente & Ventajas & Riesgos \\
\midrule
Presentación de lanzamiento & Narrativa concentrada; facilita crear hitos & Cansancio de los usuarios; difícil demostrar valor a largo plazo. \\
Benchmark & Comparable, difundible, útil para el juicio técnico & Propenso al sobreajuste a las clasificaciones; distante de las tareas reales. \\
Uso real & Retroalimentación natural de desarrolladores, clientes y productos & Retroalimentación lenta e incontrolable, pero la más cercana al valor comercial. \\
\bottomrule
\end{tabularx}
\end{knowledgebox}

\subsection{Ser usado también es una forma de evaluación}

Cuando los desarrolladores integran un modelo, cuando un producto lo usa para reemplazar a su modelo anterior, o cuando otros equipos lo destilan o lo reempaquetan, deja de ser solo un modelo de un paper o de una clasificación y entra en un ecosistema real. Esto tiene ventajas y también complicaciones: la ventaja es que tareas reales validan la capacidad del modelo; la complicación es que la captura de valor se vuelve más compleja y la influencia del código abierto no necesariamente se convierte de forma directa en ingresos.

\begin{importantbox}{La lógica de evaluación del desempeño real}
Una empresa de modelos no puede preguntarse solo «¿en qué lugar estamos en la clasificación?»; también debe preguntarse «¿quién lo usa, por qué lo usa, a quién reemplazó, si el costo es menor y si los clientes están dispuestos a seguir pagando?». Esa es la demostración de capacidad más importante en la etapa de empresa que cotiza en bolsa.
\end{importantbox}

\subsection{Resumen de la sección}

El centro de los lanzamientos de modelos se está desplazando de «yo digo que soy muy bueno» a «qué lograron otros usando mi modelo». Esta es también la nueva relación entre el ecosistema de código abierto, la comercialización y la capacidad de los modelos.
\section{IPO: la primera empresa de grandes modelos del mundo en cotizar en bolsa y una nueva etapa de gobierno corporativo}

La sección sobre la salida a bolsa no es solo una historia del mercado de capitales. Zhang Peng (张鹏) habla de los umbrales de tamaño de una empresa: unas decenas de personas, doscientas, quinientas y la empresa que cotiza en bolsa corresponden a problemas de gestión distintos. Con unas decenas de personas, primero hay que generar ingresos y construir la confianza del equipo; con cien o doscientas, comienza la división del trabajo y sube el costo de comunicarse y alinearse; a partir de varios cientos aparecen niveles jerárquicos y la información recorre caminos más largos; después de la salida a bolsa, las exigencias de cumplimiento normativo y de gobierno corporativo son mayores.

\begin{figure}[H]
\centering
\includegraphics[width=0.96\textwidth]{company-stage-thresholds.png}
\caption{Umbrales de etapa de una empresa emergente: 50 personas, 200 personas, 500 personas y el gobierno corporativo de una empresa que cotiza en bolsa.}
\end{figure}

\begin{knowledgebox}{Lectura de la figura: detrás de las cifras hay etapas de organización}
Las cifras de la figura no son umbrales exactos, sino saltos en la complejidad de la organización. Cuanto más grande es el equipo, menos puede depender de que la visión personal del fundador abarque todo; tiene que apoyarse en mecanismos, comités, procesos, cumplimiento normativo y mandos medios.
\end{knowledgebox}

\subsection{Ventajas e inercias del emprendimiento de origen académico}

El conductor pregunta si «de origen académico» es una descripción acertada. Zhang Peng reconoce que es bastante precisa, porque el equipo salió de la universidad. La ventaja del perfil académico es que valora la tecnología, la investigación y los problemas de largo plazo; su inercia es que puede descuidar la comercialización. Sin embargo, desde sus inicios en el laboratorio, Zhipu (智谱) ya trataba con clientes, generaba ingresos y entregaba sistemas, por lo que aprendió relativamente pronto el lenguaje de la comercialización.

\begin{importantbox}{La corrección clave del emprendimiento de origen académico}
El idealismo tecnológico de origen académico tiene que complementarse con capacidad de comercialización. Investigación, ingeniería, clientes, cotizaciones, costos y entregas deben quedar en la misma tabla.
\end{importantbox}

\subsection{Resumen de la sección}

La IPO es una nueva etapa para Zhipu. No solo implica financiamiento y notoriedad, sino también una mayor complejidad en el gobierno corporativo, el cumplimiento normativo, el flujo de la información y la disciplina comercial.
\section{Después de la salida a bolsa: disciplina financiera, inversión en investigación y desarrollo, y equilibrio comercial}

Después de salir a bolsa, los problemas de una empresa de modelos se vuelven más concretos: la inversión en investigación y desarrollo es muy alta, los ingresos comerciales deben crecer, el mercado de capitales quiere ver una ruta, los clientes quieren ver entregas estables y, al mismo tiempo, el equipo de investigación no puede perder la visión de largo plazo. En la entrevista, Zhang Peng (张鹏) no hizo promesas exageradas; en cambio, enfatizó que los ingresos en la nube, los ingresos 2B, la optimización de costos y la ruta hacia el break even mejorarán de forma gradual.

\subsection{Por qué «solo ganar dinero» no basta}

El capítulo anterior trató la disciplina financiera en la etapa de empresa cotizada; esta sección completa la otra mitad: la disciplina no puede aplastar el ideal técnico hasta convertirlo en un objetivo puramente de utilidades. De lo contrario, la empresa de modelos perdería su identidad.

El entrevistador preguntó si él estaría satisfecho si, en los próximos cinco años, Zhipu (智谱) solo ganara dinero, sin producción técnica ni contribución a la industria. La respuesta de Zhang Peng fue negativa. Esta respuesta muestra que Zhipu no considera la salida a bolsa como un punto final ni toma las utilidades como su único objetivo. Para una empresa de modelos, las utilidades son importantes, desde luego; pero si provienen de subcontratación de corto plazo, de integración o de entregas de bajo contenido técnico, mientras la capacidad del modelo deja de avanzar, la empresa pierde su identidad como empresa de grandes modelos.

\begin{warningbox}{La presión de ser empresa cotizada es un arma de doble filo}
Salir a bolsa aporta transparencia, capacidad de financiamiento y normas de gobierno corporativo, pero también trae presión financiera de corto plazo. Si una empresa de modelos se deja arrastrar por completo por las utilidades de corto plazo, puede debilitar su investigación y desarrollo de largo plazo; si solo habla de una visión de largo plazo sin demostrar su comercialización, tampoco puede sostener una inversión continua.
\end{warningbox}

\subsection{Cómo la comercialización retroalimenta la investigación}

La lógica P2P de los primeros años de Zhipu ya contenía un ciclo cerrado: la práctica de ingeniería retroalimenta la investigación. Después de la salida a bolsa, este ciclo se vuelve todavía más importante. Las necesidades de los clientes revelan las carencias reales del modelo, las llamadas en la nube revelan los problemas de costo y latencia, y los proyectos 2B revelan problemas de seguridad, permisos, integración y confiabilidad. Si estos problemas entran en la cola de investigación y desarrollo, acercan el modelo al mundo real; si solo se tratan como una carga de posventa y entrega, separan la investigación del negocio.

\begin{importantbox}{El ciclo correcto después de la salida a bolsa}
La inversión en investigación y desarrollo produce capacidad del modelo; la capacidad del modelo produce valor para el cliente; el valor para el cliente produce ingresos y retroalimentación real, y esa retroalimentación real vuelve a alimentar la investigación y el desarrollo. Este ciclo es más sostenible que limitarse a «recaudar fondos y gastarlos en entrenar modelos».
\end{importantbox}

\subsection{Resumen de la sección}

Tras su salida a bolsa, Zhipu debe demostrar dos cosas a la vez: que puede seguir persiguiendo la AGI y que puede sostener esa búsqueda con ingresos comerciales y disciplina de gobierno corporativo. Ambas cosas no se oponen por naturaleza, pero exigen un equilibrio organizacional muy sólido.
\section{El CEO como puente: investigación, comercialización y expectativas de la organización}

Antes y después de la salida a bolsa, el papel del CEO cambia. Zhang Peng (张鹏) dice que ahora dedica más energía a la operación diaria de la empresa, en especial al área de atención al cliente, la orientación al mercado y la comercialización. También menciona que el CEO muchas veces es un puente, quien monta el escenario: hace que los departamentos de investigación y de tecnología y los equipos de comercialización puedan comunicarse, y que todos desplieguen su imaginación, su capacidad y su capacidad de ejecución.

\subsection{La investigación y la comercialización no pueden avanzar desconectadas}

En una empresa de modelos es fácil que haya dos piezas desconectadas: el equipo de investigación se ocupa de las capacidades del modelo y el equipo comercial, de las entregas al cliente. Zhang Peng subraya que internamente hay que alinear de forma continua los cambios del mercado, la investigación y el desarrollo tecnológicos y las necesidades de comercialización. De lo contrario, la investigación puede producir algo vistoso pero difícil de llevar a la práctica, y el área comercial puede sacrificar la dirección tecnológica de largo plazo a cambio de ingresos de corto plazo.

\begin{importantbox}{La doble tarea del CEO de una empresa de modelos}
Por un lado, debe proteger la investigación de largo plazo y las capacidades del modelo; por otro, debe asegurar el valor para el cliente, los ingresos, el cumplimiento normativo y la eficiencia de la organización. Lo central del CEO no es decidir por todos, sino lograr que estos sistemas puedan escucharse entre sí.
\end{importantbox}

\subsection{Algunos momentos conmovedores}

El papel de puente del CEO a veces resulta abstracto, por eso esta sección lo aterriza con algunos momentos concretos. Muestran que una empresa de modelos no se demuestra solo en artículos y presentaciones de lanzamiento.

Al final de la entrevista, Zhang Peng recuerda algunos momentos. Uno es haber pasado casi medio año en Shenzhen por un cliente grande y, al final, no volver a Pekín con las manos vacías, sino con contratos por decenas de millones; esto muestra que los clientes están dispuestos a pagar por la tecnología. Otro es haber repartido sobres rojos a los asistentes con un agente inteligente durante una presentación de lanzamiento: aunque el monto se escribió con un dígito equivocado, los sobres rojos sí se enviaron, y se comentó que era el primer sobre rojo que la IA le daba a la humanidad. Otro más es que, después de que GLM-4.5 se publicó como código abierto, desarrolladores y productos del extranjero empezaron a usarlo, e incluso aparecieron toda clase de reutilizaciones, destilaciones y reempaquetados.

Estos momentos muestran en conjunto una misma cosa: la satisfacción de una empresa de modelos no viene solo de los benchmark, sino también de los contratos con clientes, el uso real, el reconocimiento de los desarrolladores y la repercusión de su tecnología.

\begin{knowledgebox}{Del benchmark a la repercusión real}
Un benchmark demuestra capacidad; un contrato con un cliente demuestra valor; el uso del código abierto demuestra ecosistema; el evento de lanzamiento demuestra experiencia. Una empresa de modelos tiene que gestionar al mismo tiempo estos distintos tipos de retroalimentación.
\end{knowledgebox}

\subsection{Resumen de la sección}

El relato de Zhang Peng como CEO no busca el «estrellato»; se parece más a un papel de puente: conecta la investigación, la ingeniería, el área comercial, los clientes, la gobernanza y las expectativas del equipo. Esta es también la otra cara de la «sensación de cemento» de Zhipu (智谱).
\section{Idealismo y comercialización: quienes abren camino}

Por último, el texto vuelve a los juicios de valor: una empresa que cotiza en bolsa debe generar ganancias, pero Zhang Peng (张鹏) sigue anteponiendo la AGI y la contribución a la industria. En este punto toda la entrevista pasa de la historia de la empresa a su ideal de largo plazo.

Al final de la entrevista, le preguntaron a Zhang Peng qué elegiría entre construir una empresa que logre la AGI y construir una empresa muy rentable. Eligió la AGI, aunque también subrayó que ambas metas no se oponen: si se logra la AGI, también se tendrá una gran empresa en términos comerciales. Además, no estaría satisfecho con una Zhipu (智谱) que solo generara ganancias sin producir tecnología ni contribuir a la industria.

\begin{figure}[H]
\centering
\includegraphics[width=0.96\textwidth]{idealism-commercialization.png}
\caption{Idealismo y comercialización: ideal tecnológico, cultura de ingeniería, valor para el cliente, gobierno corporativo e identidad de pionero.}
\end{figure}

\begin{knowledgebox}{Lectura de la figura: Zhipu busca un camino equilibrado}
La figura va del ideal tecnológico al valor para el cliente y de ahí al gobierno corporativo tras la salida a bolsa, para volver al final a la idea de «pionero». Zhang Peng no insiste en proclamar la AGI sin más ni en ser solo una máquina de generar ganancias, sino en convertir «hacer que las máquinas piensen como los seres humanos» en un sistema de ingeniería capaz de servir a los clientes y a la sociedad.
\end{knowledgebox}

\subsection{Hacer que las máquinas piensen como los seres humanos}

El idealismo no puede quedarse en la palabra «AGI». Esta sección descompone el slogan de Zhipu para ver cómo apunta a la vez a un objetivo tecnológico y a un valor social.

Zhang Peng explica cómo entiende Zhipu la AGI a partir del slogan de la empresa: hacer que las máquinas piensen como los seres humanos y, a su vez, que fortalezcan a la sociedad humana. La frase tiene dos partes: la primera es un objetivo tecnológico, que las máquinas tengan mayor capacidad de comprensión, razonamiento y acción; la segunda es un objetivo de valor, que la tecnología genere un valor real en lugar de quedarse en el laboratorio.

\begin{importantbox}{Cómo volver operativa la AGI}
Si las definiciones de AGI difieren entre sí, para juzgar si una empresa persigue la AGI no basta con escuchar sus consignas: hay que ver en qué tareas por etapas, rutas tecnológicas, trayectorias de producto y valor para el cliente descompone su objetivo de largo plazo.
\end{importantbox}

\subsection{Pioneros, no solo proyectos estrella}

Lo anterior trata de «qué hacer»; aquí se trata de «cómo se quiere ser recordado». La palabra «pionero» es más sobria que «empresa estrella» y encaja mejor con el relato que Zhipu hace de sí misma.

Al inicio del programa se mencionó que hay quien compara a Zhipu con el «cemento»: hace un trabajo sólido, pero no despierta mucho entusiasmo. Al final, Zhang Peng dijo que espera que la historia recuerde a Zhipu como pionera de la AGI, como quien abrió camino. Este relato no contradice la imagen del «cemento»: quien abre camino no necesariamente es quien más atención atrae, pero tiene que pavimentar el camino durante mucho tiempo, tropezar con los obstáculos, convertir la investigación en aplicaciones, servir a los clientes y asumir las exigencias del gobierno corporativo.

\subsection{Resumen de la sección}

Zhipu se posiciona en un camino equilibrado entre el idealismo y la comercialización. La salida a bolsa no es el punto final, sino lo que permite a una empresa que persigue la AGI seguir avanzando bajo exigencias de gobierno corporativo más estrictas.
\section{Términos clave: índice de conceptos de este episodio}

Las secciones anteriores abarcan tecnología, marco institucional, organización y mercados de capitales. Aquí se reúnen los conceptos principales para facilitar las referencias cruzadas con los reportes trimestrales sobre modelos grandes, las revisiones sobre Agent y las entrevistas sobre producto.

\begin{longtable}{p{0.24\textwidth}p{0.32\textwidth}p{0.35\textwidth}}
\toprule
Término & Explicación breve & Papel en este episodio \\
\midrule
Inteligencia cognitiva & Etapa de comprensión, razonamiento y conocimiento posterior a la inteligencia perceptiva & Punto de referencia de la orientación inicial de Zhipu (智谱). \\
P2P & Paper to Project / Product & Tradición de transferencia de resultados de Tsinghua (清华). \\
Transferencia de resultados científicos y tecnológicos & Convertir, conforme a la normativa, los resultados de investigación en activos y productos de una empresa & Vía institucional por la que se fundó Zhipu. \\
GPT-3 & Impacto externo temprano del paradigma de escala & Cambió la visión técnica de los modelos grandes. \\
ChatGPT & Momento en que los modelos grandes irrumpieron en la percepción de los usuarios & Aceleró de golpe el ritmo de la industria. \\
Scaling Law & Ley según la cual la capacidad del modelo crece con la escala y la retroalimentación & Paradigma de fondo de la competencia entre modelos grandes. \\
Post-training & Entrenamiento posterior al preentrenamiento con preferencias, tareas, herramientas y retroalimentación & Etapa decisiva para que el modelo pase de «saber hablar» a «ser utilizable». \\
RL & Reinforcement Learning, aprendizaje por refuerzo & Zhang Peng (张鹏) considera que traerá una nueva forma de cómputo y un nuevo eje de crecimiento de capacidades. \\
Agent & Agente que invoca herramientas y ejecuta tareas & Conecta la capacidad del modelo con aplicaciones reales. \\
DeepSeek & Caso de la industria que representa el código abierto, la eficiencia y el impacto de ingeniería & Llevó a las empresas de modelos a replantear sus lanzamientos y el código abierto. \\
Código abierto / código cerrado & Combinación estratégica entre difusión del ecosistema y control comercial & Las empresas de modelos deben elegir entre ambos. \\
Benchmark & Evaluación que compara la capacidad de los modelos con tareas uniformes & Ayuda a comparar, pero no sustituye el uso real. \\
Uso real & Adopción continua del modelo por desarrolladores, clientes y productos & Prueba de capacidad más relevante en la etapa de empresa cotizada. \\
IPO & Salida a bolsa: entrada al mercado público y a una gobernanza de cumplimiento más estricta & Nueva etapa de Zhipu. \\
Pionero & Quien invierte temprano, abre camino temprano y asume temprano la incertidumbre & Nota histórica que Zhang Peng espera que deje Zhipu. \\
\bottomrule
\end{longtable}

\subsection{Resumen de la sección}

En conjunto, estos conceptos describen la cadena completa de una empresa de modelos: del laboratorio, la transferencia de resultados y el impacto de los paradigmas técnicos, al ecosistema de código abierto, la comercialización, la gobernanza como empresa cotizada y la narrativa de largo plazo sobre la AGI.
\section{Síntesis y ampliación}

\subsection{Conclusiones principales}

\begin{enumerate}
\item El punto de partida de Zhipu (智谱) es la tradición P2P del Laboratorio de Ingeniería del Conocimiento de Tsinghua, y no solo el auge de los modelos grandes.
\item La inteligencia cognitiva fue la respuesta que Zhipu dio en su etapa temprana a la pregunta «¿qué viene después de la inteligencia perceptiva?».
\item GPT-3, ChatGPT y DeepSeek sacudieron, respectivamente, la imaginación técnica, la percepción de los usuarios y el discurso sobre el código abierto y la eficiencia.
\item La Scaling Law está pasando de ser solo una ley de escala del preentrenamiento a extenderse a la multimodalidad, el posentrenamiento, el RL y los Agent.
\item El código abierto y el código cerrado no forman una dicotomía moral, sino una combinación estratégica de ecosistema, negocio y límites de seguridad.
\item La visibilidad de las empresas de modelos se está desplazando del discurso de los eventos de lanzamiento hacia los resultados de uso reales; la adopción sostenida por parte de desarrolladores y clientes será cada vez más importante.
\item La IPO marca la entrada de la empresa en una nueva etapa de gobernanza, cumplimiento normativo y complejidad organizacional.
\item La ventaja de emprender desde la academia es un ideal técnico de largo plazo, pero debe complementarse con la comercialización y el valor para el cliente.
\item Zhipu se define a sí misma como pionera de la AGI: necesita resultados técnicos, pero también implementación de ingeniería y comercialización.
\end{enumerate}

\subsection{Preguntas abiertas}

Se dejan preguntas abiertas al final porque, aunque Zhipu ya cotiza en bolsa, la forma de largo plazo de las empresas de modelos grandes todavía no se ha definido. Las siguientes preguntas determinarán la competencia real después de «la primera acción de modelos grandes del mundo».

\begin{itemize}
\item Una vez en bolsa, ¿cómo mantiene una empresa de modelos grandes el equilibrio entre la presión del mercado de capitales y la inversión de largo plazo en AGI?
\item ¿Puede la influencia de los modelos de código abierto en el ecosistema convertirse en un valor comercial estable?
\item ¿Se convertirán los Agent y el RL en los principales ejes de crecimiento de la siguiente etapa de la Scaling Law?
\item ¿Las empresas chinas de modelos grandes se orientarán a ser empresas de aplicaciones, empresas 2B, plataformas de código abierto o una combinación de varias rutas?
\item ¿Cómo se convierte «hacer que las máquinas piensen como las personas» en metas por etapas que puedan evaluarse, entregarse y regularse?
\end{itemize}

\subsection{Lecturas adicionales}

\begin{itemize}
\item EP136, informe trimestral global de modelos grandes: para entender la diferenciación entre empresas de modelos, el Coding y el modelo OS.
\item EP127, conversación de fin de año del informe trimestral de modelos grandes: para entender la AI War, las alianzas y el Online Learning.
\item EP139, historia técnica de los Agent: para entender por qué los Agent se volvieron la vía de implementación de los modelos.
\item EP117, historia de los paradigmas de modelos, la Infra y los datos: para entender la trayectoria histórica de los modelos grandes.
\item Informes técnicos de modelos de código abierto como DeepSeek, GLM, Qwen y Kimi: para comparar las estrategias de código abierto y la evolución de sus capacidades.
\end{itemize}

\begin{importantbox}{Valoración final}
Lo que más vale la pena llevarse de esta entrevista con Zhipu es el significado detrás de la palabra «pionero»: no es quien mejor sabe generar visibilidad, sino quien abre el camino primero cuando lo institucional, lo técnico, la ingeniería y el negocio son todavía inciertos.
\end{importantbox}

\end{document}
